\documentclass[11pt]{article}
\usepackage[T1]{fontenc}
\usepackage[utf8]{inputenc}
\usepackage{lmodern}
\usepackage[english]{babel}
\usepackage{amsmath,amssymb,mathtools,array,booktabs,pdflscape,adjustbox,diagbox}
\usepackage{geometry}
\usepackage{microtype}
\geometry{a4paper,margin=1.45cm}
\setlength{\parindent}{1.25em}
\setlength{\parskip}{0pt}
\makeatletter
\renewcommand\footnotesize{\@setfontsize\footnotesize{7.5}{8.5}\selectfont}
\makeatother
\providecommand{\widebar}[1]{\overline{#1}}
\providecommand{\A}{\Delta}
\providecommand{\D}{\Delta}
\allowdisplaybreaks
\setlength{\emergencystretch}{10em}
\sloppy
\hbadness=10000
\vbadness=10000
\hfuzz=2pt
\vfuzz=10pt
\raggedbottom
\providecommand{\R}{\varrho}
\providecommand{\hata}{\widehat{a}}
\providecommand{\hatv}{\widehat{v}}
\providecommand{\modu}[1]{\;\operatorname{mod} #1}
\providecommand{\mcell}[1]{\ensuremath{\begin{gathered}#1\end{gathered}}}
\providecommand{\pair}[2]{(#1\,|\,#2)}
\providecommand{\eps}{\varepsilon}
\providecommand{\rhoidx}{\varrho}

% Batch 20 macros
\providecommand{\dd}{\mathrm d}
\providecommand{\Div}{\operatorname{Div}}
\providecommand{\G}{\mathfrak G}
\providecommand{\bd}{\bar{\delta}}
\providecommand{\p}{\partial}


% Batch 21 ideal/function-field macros
\providecommand{\frS}{\mathfrak S}
\providecommand{\frp}{\mathfrak p}
\providecommand{\fra}{\mathfrak a}
\providecommand{\frb}{\mathfrak b}
\providecommand{\frc}{\mathfrak c}
\providecommand{\frf}{\mathfrak f}
\providecommand{\frr}{\mathfrak r}
\providecommand{\frm}{\mathfrak m}
\providecommand{\frD}{\mathfrak D}
\providecommand{\calA}{\mathcal A}
\providecommand{\calB}{\mathcal B}
\providecommand{\calN}{\mathcal N}


\providecommand{\p}{\partial}
\providecommand{\dd}{\mathrm d}
\providecommand{\modu}[1]{\;\operatorname{mod} #1}
\providecommand{\mM}{\mathfrak{M}}
\providecommand{\mN}{\mathfrak{N}}
\providecommand{\mL}{\mathfrak{L}}
\providecommand{\mG}{\mathfrak{G}}
\providecommand{\mH}{\mathfrak{H}}
\providecommand{\mH}{\mathfrak{H}}
\providecommand{\mA}{\mathfrak{A}}
\providecommand{\mB}{\mathfrak{B}}
\providecommand{\mX}{\mathfrak{X}}
\providecommand{\mY}{\mathfrak{Y}}
\providecommand{\ma}{\mathfrak{a}}
\providecommand{\mb}{\mathfrak{b}}
\providecommand{\mx}{\mathfrak{x}}
\providecommand{\my}{\mathfrak{y}}
\providecommand{\mz}{\mathfrak{z}}
\providecommand{\me}{\mathfrak{e}}
\providecommand{\mbarA}{\overline{\mathfrak{A}}}
\providecommand{\mbarB}{\overline{\mathfrak{B}}}
\providecommand{\mbar}[1]{\overline{#1}}
\providecommand{\mC}{\mathfrak{C}}
\providecommand{\mE}{\mathfrak{E}}
\providecommand{\mF}{\mathfrak{F}}
\providecommand{\mT}{\mathfrak{T}}
\providecommand{\mW}{\mathfrak{W}}
\providecommand{\md}{\mathfrak{d}}
\providecommand{\mf}{\mathfrak{f}}
\providecommand{\mm}{\mathfrak{m}}
\providecommand{\mn}{\mathfrak{n}}
\providecommand{\mo}{\mathfrak{o}}
\providecommand{\mpideal}{\mathfrak{p}}
\providecommand{\mq}{\mathfrak{q}}
\providecommand{\mt}{\mathfrak{t}}
\providecommand{\mv}{\mathfrak{v}}
\providecommand{\bara}{\overline{\mathfrak{a}}}



% Batch 25 module/residue macros
\providecommand{\mU}{\mathfrak{U}}
\providecommand{\mP}{\mathfrak{P}}
\providecommand{\mQ}{\mathfrak{Q}}
\providecommand{\mS}{\mathfrak{S}}
\providecommand{\mD}{\mathfrak{D}}
\providecommand{\mR}{\mathfrak{R}}
\providecommand{\mV}{\mathfrak{V}}
\providecommand{\ideal}[1]{\mathfrak{#1}}

% Batch 31 polynomial-ideal/resultant macros
\providecommand{\Amod}{\mathfrak A}
\providecommand{\Bmod}{\mathfrak B}
\providecommand{\Gmod}{\mathfrak G}
\providecommand{\Rmod}{\mathfrak R}
\providecommand{\Smod}{\mathfrak S}
\providecommand{\mideal}{\mathfrak m}
\providecommand{\nideal}{\mathfrak n}
\providecommand{\gideal}{\mathfrak g}
\providecommand{\bb}{\mathfrak b}
\providecommand{\cc}{\mathfrak c}
\providecommand{\zero}[1]{\equiv 0(#1)}
\providecommand{\Norm}{N}


\providecommand{\Amod}{\mathfrak A}
\providecommand{\Bmod}{\mathfrak B}
\providecommand{\Gmod}{\mathfrak G}
\providecommand{\Mmod}{\mathfrak M}
\providecommand{\Nmod}{\mathfrak N}
\providecommand{\Smod}{\mathfrak S}
\providecommand{\mideal}{\mathfrak m}
\providecommand{\nideal}{\mathfrak n}
\providecommand{\gideal}{\mathfrak g}
\providecommand{\hideal}{\mathfrak h}
\providecommand{\jideal}{\mathfrak j}
\providecommand{\tideal}{\mathfrak t}
\providecommand{\aideal}{\mathfrak a}
\providecommand{\bideal}{\mathfrak b}
\providecommand{\zero}[1]{\equiv 0\pmod{#1}}
\providecommand{\Norm}{N}


\providecommand{\frakS}{\mathfrak S}
\providecommand{\frakM}{\mathfrak M}
\providecommand{\frakG}{\mathfrak G}
\providecommand{\frakm}{\mathfrak m}
\providecommand{\frakn}{\mathfrak n}
\providecommand{\frakp}{\mathfrak p}
\providecommand{\frakq}{\mathfrak q}
\providecommand{\frakr}{\mathfrak r}
\providecommand{\fraka}{\mathfrak a}
\providecommand{\frakb}{\mathfrak b}
\providecommand{\frakc}{\mathfrak c}
\providecommand{\frakd}{\mathfrak d}
\providecommand{\frake}{\mathfrak e}
\providecommand{\frakg}{\mathfrak g}
\providecommand{\fraks}{\mathfrak s}
\providecommand{\frako}{\mathfrak o}
\providecommand{\frK}{\mathfrak K}
\providecommand{\barfrakm}{\overline{\mathfrak m}}

\providecommand{\frakh}{\mathfrak h}
\providecommand{\frakt}{\mathfrak t}

% Batch 36 abstract ideal theory / invariant theory macros
\providecommand{\frakR}{\mathfrak R}
\providecommand{\frakT}{\mathfrak T}
\providecommand{\frakH}{\mathfrak H}
\providecommand{\frakP}{\mathfrak P}
\providecommand{\frakQ}{\mathfrak Q}
\providecommand{\frakZ}{\mathfrak Z}
\providecommand{\frakN}{\mathfrak N}
\providecommand{\frakB}{\mathfrak B}
\providecommand{\frakF}{\mathfrak F}
\providecommand{\frakL}{\mathfrak L}
\providecommand{\frakW}{\mathfrak W}
\providecommand{\frakGg}{\mathfrak G}
\providecommand{\Endl}{\operatorname{End}}
\providecommand{\Gal}{\operatorname{Gal}}
\providecommand{\Trdeg}{\operatorname{trdeg}}


% Batch 38 continuation macros
\providecommand{\frR}{\mathfrak R}
\providecommand{\frK}{\mathfrak K}
\providecommand{\frL}{\mathfrak L}
\providecommand{\frM}{\mathfrak M}
\providecommand{\frF}{\mathfrak F}
\providecommand{\frT}{\mathfrak T}
\providecommand{\frN}{\mathfrak N}
\providecommand{\frO}{\mathfrak O}
\providecommand{\frB}{\mathfrak B}
\providecommand{\frC}{\mathfrak C}
\providecommand{\frA}{\mathfrak A}


% Batch 41 Noether macros
\providecommand{\mR}{\mathfrak{R}}
\providecommand{\bR}{\overline{\mathfrak{R}}}
\providecommand{\mK}{\mathfrak{K}}
\providecommand{\mL}{\mathfrak{L}}
\providecommand{\mS}{\mathfrak{S}}
\providecommand{\mT}{\mathfrak{T}}
\providecommand{\mZ}{\mathfrak{Z}}
\providecommand{\mO}{\mathfrak{O}}
\providecommand{\mo}{\mathfrak{o}}
\providecommand{\mD}{\mathfrak{D}}
\providecommand{\mE}{\mathfrak{E}}
\providecommand{\mP}{\mathfrak{P}}
\providecommand{\mf}{\mathfrak{f}}
\providecommand{\mpideal}{\mathfrak{p}}
\providecommand{\mq}{\mathfrak{q}}
\providecommand{\mt}{\mathfrak{t}}
\providecommand{\mOmega}{\mathfrak{\Omega}}
\providecommand{\bP}{\overline{P}}
\providecommand{\bq}{\overline{\mathfrak{q}}}
\providecommand{\bp}{\overline{\mathfrak{p}}}
\providecommand{\ba}{\overline{\mathfrak{a}}}
\providecommand{\bb}{\overline{\mathfrak{b}}}
\providecommand{\tr}{\operatorname{Tr}}
\providecommand{\Norm}{\operatorname{N}}


\providecommand{\rank}{\operatorname{rank}}

% Batch 46 macros
\providecommand{\frakA}{\mathfrak A}
\providecommand{\frakB}{\mathfrak B}
\providecommand{\frakM}{\mathfrak M}
\providecommand{\frakN}{\mathfrak N}
\providecommand{\frako}{\mathfrak o}
\providecommand{\fraka}{\mathfrak a}
\providecommand{\frakc}{\mathfrak c}
\providecommand{\frakp}{\mathfrak p}
\providecommand{\frakO}{\mathfrak O}
\providecommand{\frakP}{\mathfrak P}
\providecommand{\frakD}{\mathfrak D}
\providecommand{\frakG}{\mathfrak G}
\providecommand{\frakK}{\mathfrak K}
\providecommand{\Sp}{\operatorname{Sp}}
\providecommand{\rank}{\operatorname{rank}}
\providecommand{\Trdeg}{\operatorname{Trdeg}}
\providecommand{\charac}{\operatorname{char}}



% Batch 47 macros
\providecommand{\Om}{\Omega}
\providecommand{\Omfrakp}{\Omega_{\mathfrak p}}
\providecommand{\frakA}{\mathfrak A}
\providecommand{\frakB}{\mathfrak B}
\providecommand{\frakC}{\mathfrak C}
\providecommand{\frakG}{\mathfrak G}
\providecommand{\frakH}{\mathfrak H}
\providecommand{\frakP}{\mathfrak P}
\providecommand{\frakp}{\mathfrak p}
\providecommand{\frakq}{\mathfrak q}
\providecommand{\frakr}{\mathfrak r}
\providecommand{\frakS}{\mathfrak S}
\providecommand{\Cl}{\operatorname{Cl}}
\providecommand{\ind}{\operatorname{ind}}
\providecommand{\N}{\operatorname{N}}
\providecommand{\Mat}{\operatorname{M}}
\providecommand{\Gal}{\operatorname{Gal}}
\providecommand{\Norm}{\operatorname{N}}
\providecommand{\nr}{\operatorname{nr}}
\providecommand{\Tr}{\operatorname{Tr}}
\providecommand{\Br}{\operatorname{Br}}


% Batch 48 macros
\providecommand{\frR}{\mathfrak R}
\providecommand{\frS}{\mathfrak S}
\providecommand{\frT}{\mathfrak T}
\providecommand{\frM}{\mathfrak M}
\providecommand{\frN}{\mathfrak N}
\providecommand{\frA}{\mathfrak A}
\providecommand{\frB}{\mathfrak B}
\providecommand{\frX}{\mathfrak X}
\providecommand{\frY}{\mathfrak Y}
\providecommand{\frZ}{\mathfrak Z}
\providecommand{\frG}{\mathfrak G}
\providecommand{\frH}{\mathfrak H}
\providecommand{\frL}{\mathfrak L}
\providecommand{\frU}{\mathfrak U}
\providecommand{\frD}{\mathfrak D}
\providecommand{\frC}{\mathfrak C}
\providecommand{\frP}{\mathfrak P}
\providecommand{\calR}{\mathcal R}
\providecommand{\calS}{\mathcal S}
\providecommand{\calA}{\mathcal A}
\providecommand{\calB}{\mathcal B}
\providecommand{\calT}{\mathcal T}
\providecommand{\End}{\operatorname{End}}
\providecommand{\Aut}{\operatorname{Aut}}
\providecommand{\Mat}{\operatorname{Mat}}
\providecommand{\rad}{\operatorname{rad}}
\providecommand{\rk}{\operatorname{rk}}
\providecommand{\id}{\operatorname{id}}
\providecommand{\op}{\operatorname{op}}


% Batch 49 macros
\providecommand{\Alg}{\mathcal A}
\providecommand{\Abar}{\overline{A}}
\providecommand{\Bbar}{\overline{B}}
\providecommand{\Kbar}{\overline{K}}
\providecommand{\Zbar}{\overline{Z}}
\providecommand{\Gcal}{\mathfrak G}
\providecommand{\Hcal}{\mathfrak H}
\providecommand{\Scal}{\mathfrak S}
\providecommand{\Tcal}{\mathfrak T}
\providecommand{\Zcal}{\mathfrak Z}
\providecommand{\Kcal}{\mathfrak K}
\providecommand{\Pcal}{\mathfrak P}
\providecommand{\rad}{\operatorname{rad}}
\providecommand{\Sp}{\operatorname{Sp}}
\providecommand{\End}{\operatorname{End}}
\providecommand{\Mat}{\operatorname{Mat}}
\providecommand{\Inv}{\operatorname{Inv}}


% Batch 50 macros
\providecommand{\Sp}{\operatorname{Sp}}
\providecommand{\Tr}{\operatorname{Tr}}
\providecommand{\Norm}{\operatorname{N}}
\providecommand{\Mat}{\operatorname{Mat}}
\providecommand{\Gal}{\operatorname{Gal}}
\providecommand{\Cl}{\operatorname{Cl}}
\providecommand{\Gg}{\mathfrak G}
\providecommand{\Hh}{\mathfrak H}
\providecommand{\Ii}{\mathfrak I}
\providecommand{\Jj}{\mathfrak J}
\providecommand{\Aa}{\mathfrak A}
\providecommand{\Bb}{\mathfrak B}
\providecommand{\Cc}{\mathfrak C}
\providecommand{\Oo}{\mathfrak O}
\providecommand{\oo}{\mathfrak o}
\providecommand{\pp}{\mathfrak p}
\providecommand{\PP}{\mathfrak P}
\providecommand{\LL}{\mathfrak L}
\providecommand{\RR}{\mathfrak R}
\providecommand{\ZZ}{\mathfrak Z}
\providecommand{\ee}{\mathfrak e}


% Batch 51 macros
\providecommand{\Sp}{\operatorname{Sp}}
\providecommand{\Trdeg}{\operatorname{Trdeg}}
\providecommand{\frO}{\mathfrak O}
\providecommand{\fro}{\mathfrak o}
\providecommand{\frh}{\mathfrak h}
\providecommand{\frB}{\mathfrak B}
\providecommand{\frA}{\mathfrak A}
\providecommand{\frM}{\mathfrak M}
\providecommand{\frD}{\mathfrak D}
\providecommand{\frd}{\mathfrak d}
\providecommand{\frK}{\mathfrak K}
\providecommand{\frg}{\mathfrak g}
\providecommand{\frp}{\mathfrak p}

% Extra macros collected from the cumulative source for standalone paper builds.
\providecommand{\mat}[1]{\begin{pmatrix}#1\end{pmatrix}}
\providecommand{\Afield}{\mathfrak A}
\providecommand{\Bfield}{\mathfrak B}
\providecommand{\Cfield}{\mathfrak C}
\providecommand{\Sdomain}{\mathfrak S}
\providecommand{\Kfield}{\mathfrak K}
\providecommand{\Rfield}{\mathfrak R}
\providecommand{\Xreal}{\mathfrak X}
\providecommand{\Yreal}{\mathfrak Y}
\providecommand{\Kfield}{\mathfrak{K}}
\providecommand{\Ssys}{\mathfrak{S}}
\providecommand{\Mfield}{\mathfrak{M}}
\providecommand{\tq}{\tau}
\providecommand{\ds}{\displaystyle}
\providecommand{\Lsys}{\mathfrak{L}}
\providecommand{\Jdom}{\mathfrak{J}}
\providecommand{\Gdom}{\mathfrak{G}}
\providecommand{\Omegaint}{[\Omega]}
\providecommand{\Hgrp}{\mathfrak{H}}
\providecommand{\GaloisRes}{\Phi}
\providecommand{\pol}[2]{\left(#1\frac{\partial}{\partial #2}\right)}
\providecommand{\Deltaop}{\Delta\!\left(\frac{\partial}{\partial x},\frac{\partial}{\partial y},\frac{\partial}{\partial z}\right)}
\providecommand{\Omegaop}{\Omega}
\providecommand{\bdelta}{\bar{\delta}}
\providecommand{\tmod}[1]{\;(\operatorname{mod} #1)}
\providecommand{\eqzero}[1]{\equiv 0\,(#1)}
\providecommand{\mJ}{\mathfrak J}
\providecommand{\mS}{\mathfrak S}
\providecommand{\mo}{\mathfrak o}
\providecommand{\ma}{\mathfrak a}
\providecommand{\mb}{\mathfrak b}
\providecommand{\dx}{\dd x}
\providecommand{\dy}{\dd y}
\providecommand{\mc}{\mathfrak{c}}
\providecommand{\OO}{\Omega}
\providecommand{\Th}{\Theta}
\providecommand{\iso}{\simeq}
\providecommand{\normlhd}{\triangleleft}
\providecommand{\mh}{\mathfrak{h}}
\providecommand{\ml}{\mathfrak{l}}
\providecommand{\mr}{\mathfrak{r}}
\providecommand{\ann}{\operatorname{Ann}}
\providecommand{\tuple}[1]{(#1)}
\begin{document}
% Generated from the final cumulative English Noether TeX for reader-facing packaging.
% Source file: Noether_FINAL_Cumulative_English_Translation_MODERN_TEX_numbered_papers_complete_AUDITED.tex
% Paper: 25; source line range: 14246-14296
\section*{25. Elimination Theory and Ideal Theory}

\begin{center}
J. Ber. d. DMV 33 (1924), pp. 116--120
\end{center}

In what follows I would like to show how elimination theory can be built up in exact parallel with the theory of zeros for polynomials in one indeterminate. This construction rests on a connection between Hentzelt's elimination theory and the methods of field theory, as developed by Steinitz, and ideal theory, as I myself have developed it.\footnote{Detailed presentation and literature references are in Math. Ann. 90 and in a forthcoming paper.}

I first want to sketch the question in the case of polynomials in one indeterminate. As coefficient domain let a fixed field $P$ be taken once for all, to which the concepts prime function, etc., refer. Here $P$ may be assumed to be abstractly defined in the sense of Steinitz, but in the special case it may naturally coincide with the field of rational numbers or the field of all complex numbers -- according to the assumption formerly usual in elimination theory. The theory of zeros for polynomials in one indeterminate now splits into four steps:

1. \emph{The decomposition theorem}, which permits the reduction of the question to one for prime functions -- polynomials irreducible in $P$. For from
\[
 a(x)=p_1(x)^{\varrho_1}\cdots p_r(x)^{\varrho_r}=q_1(x)\cdots q_r(x)
\]
-- where the $p(x)$ denote distinct prime functions and the $q(x)$ therefore denote ``largest primary functions'' -- it follows that $a(x)$ vanishes at all and only the zeros of the prime functions $p(x)$.

2. \emph{The construction of the zero field for a prime function $p(x)$.} Such a zero field $\mathfrak R=P(\alpha)$ -- where $\alpha$ denotes a zero of $p(x)$ -- can always, by Steinitz, be constructed as an extension field of $P$ isomorphic to the residue-class field $(\mathfrak R)$ of $p(x)$. Conversely, every zero field lying in an extension field of $P$ is isomorphic to the residue-class field $(\mathfrak R)$; zero field and residue-class field are thus abstractly identical.

3. \emph{The decomposition of $p(x)$ into linear factors in the Galois field determined by $p(x)$.} By repeating the construction given under 2. one arrives at the essentially uniquely determined Galois field $P(\alpha_1,\ldots,\alpha_m)$, which is precisely sufficient for $p(x)$ to split into linear factors.

4. \emph{The meaning of multiplicity for the original polynomial $a(x)$.} By 1., 2., 3., the question of the zeros, and the corresponding decomposition into linear factors, is settled for an arbitrary polynomial $a(x)$. The multiplicity may be regarded as the degree of the individual primary function $q(x)$ -- that is, the number of residue classes modulo $q(x)$ linearly independent with respect to $P$; for in the case of the algebraically closed coefficient domain $P$, where the $p(x)$ become linear, this number is precisely the multiplicity of the zeros.

I now pass to general elimination theory, and therefore take as basis the domain $\mathfrak S$ of all polynomials in $x_1,\ldots,x_n$ with coefficients from $P$; elimination theory can then be defined as the \emph{theory of the zeros of an ideal} from this polynomial domain. By an ideal $\mathfrak a$ I understand, as usual, a system of polynomials such that together with $f$ and $g$ also $f-g$, and together with $f$ also $Af$, belongs to $\mathfrak a$, where $A$ is understood to be an arbitrary polynomial from $\mathfrak S$. By a zero of $\mathfrak a$ I understand every value system $\alpha_1,\ldots,\alpha_n$, belonging to an extension field of $P$, such that $f(\alpha)$ vanishes for every polynomial $f$ belonging to the ideal. If $f_1,\ldots,f_t$ denotes an always existing ideal basis of $\mathfrak a$ -- so that $f=A_1f_1+\cdots+A_tf_t$ holds for every $f$ from $\mathfrak a$ -- then the zeros of $\mathfrak a$ are plainly identical with the zeros of $f_1,\ldots,f_t$. Since, conversely, every system of finitely many polynomials can be regarded as a basis of the ideal derived from it, the definition of zeros just given for elimination theory is identical with the customary one, but avoids the difficulty lying in the arbitrary singling out of a basis. For polynomials in one indeterminate, where only principal ideals are involved, this difficulty does not arise, since the basis polynomial is uniquely determined up to a unit -- a quantity from $P$ as factor -- but even there the ideal-theoretic conception will give a more exact insight.

Elimination theory now splits, correspondingly to the case of one indeterminate, into four steps:

I. \emph{The decomposition theorem} -- the representation of an ideal by largest primary components -- which permits the reduction of the question to one for prime ideals. Here an ideal is called a prime ideal if it divides a product only when it divides at least one factor; it is called a primary ideal if it divides at least one factor or a power of each factor. To every primary ideal $\mathfrak q$ there belongs one and only one associated prime ideal $\mathfrak p$, which is a divisor of $\mathfrak q$ and a power of which is divisible by $\mathfrak q$. One has the representation as least common multiple:
\[
 \mathfrak a=[\mathfrak q_1,\ldots,\mathfrak q_r];\qquad
 \mathfrak p_1^{\sigma_1}\cdots \mathfrak p_r^{\sigma_r}\equiv 0\pmod{\mathfrak a}.
\]
Here the $\mathfrak q$ are largest primary components; that is, each $\mathfrak q$ is primary, but the least common multiple of any two $\mathfrak q$'s is not primary. Moreover, without loss of generality it may be assumed that no $\mathfrak q$ divides the least common multiple of the remaining ones, so that no $\mathfrak q$ can simply be omitted; the $\mathfrak p$ are the associated prime ideals of the $\mathfrak q$. Now it is not the $\mathfrak q$, but -- and herein lies the analogue to the case of one indeterminate -- their associated prime ideals $\mathfrak p$ that are uniquely determined by $\mathfrak a$; and since each $\mathfrak p$ is a divisor of $\mathfrak a$, and conversely a product of powers of the $\mathfrak p$ is divisible by $\mathfrak a$, the zeros of the ideal therefore coincide with those of its associated prime ideals.

II. \emph{The construction of the zero field for a prime ideal $\mathfrak p$.} By definition the residue classes modulo a prime ideal form a ring without zero divisors, containing at least one element distinct from the zero element as soon as $\mathfrak p$ is different from the unit ideal; by forming quotients this ring can therefore be extended to the residue-class field $(\mathfrak R)$ of $\mathfrak p$. One can always construct an extension field $\mathfrak R$ of $P$ isomorphic to $(\mathfrak R)$ which becomes the zero field of $\mathfrak p$; that is, $\mathfrak R=P(\alpha_1,\ldots,\alpha_n)$, where $\alpha_1,\ldots,\alpha_n$ denotes a zero of $\mathfrak p$. Here $\mathfrak R$ has transcendence degree $k$ $(0\leq k<n)$ over $P$, hence contains $k$ elements algebraically independent over $P$, whereas any $(k+1)$ elements become algebraically dependent over $P$. Conversely, every zero field of transcendence degree $k$ lying in an extension field of $P$ is isomorphic to the residue-class field $(\mathfrak R)$; zero field of transcendence degree $k$ and residue-class field are therefore abstractly identical.

III. \emph{Decomposition of the norm of $\mathfrak p$ into linear factors in the Galois field determined by $\mathfrak p$.} If the indeterminates $x_1,\ldots,x_n$ are regarded as generated from original indeterminates $y_1,\ldots,y_n$ by a linear transformation with indeterminates as coefficients, then the zero field $\mathfrak R$ of transcendence degree $k$ can be regarded as a finite algebraic extension field of $P(x_{i+1},\ldots,x_n)$ $(k=n-i)$. One can therefore pass to the Galois field $\mathfrak R$ over $P(x_{i+1},\ldots,x_n)$, in which the prime function $P^{(i)}(x_i,x_{i+1},\ldots,x_n)$ with the zero $\alpha_i$ splits into linear factors. This prime function $P^{(i)}$, when one returns to the original indeterminates $y$, plays the role of the fundamental equation; for $\alpha_i$ becomes equal to the fundamental form $u_1\beta_1+\cdots+u_n\beta_n$, where $\beta$ denotes a zero system in the $y$ algebraically dependent on the parameters $x_{i+1},\ldots,x_n$. The decomposition into linear factors thus gives the product over all zeros. But the prime function $P^{(i)}$, or a power of it, may also be regarded as the norm $N(\mathfrak p)$ of $\mathfrak p$, by considering $\mathfrak p$ as a module of linear forms in the power products of $x_1,\ldots,x_{i-1}$, taking $P(x_{i+1},\ldots,x_n)$ instead of $P$ as coefficient domain, and defining the norm as the product over all elementary divisors -- of which only finitely many are distinct from the unit. For a perfect field $P$ as coefficient domain one always obtains $N(\mathfrak p)=P^{(i)}$; for an imperfect field a power of $P^{(i)}$ may occur as the norm.

IV. \emph{The norm of the original ideal and the meaning of multiplicity.} If correspondingly one regards the ideal $\mathfrak a$ successively as a module of linear forms in the power products of $x_1,\ldots,x_{i-1}$ $(i=1,2,\ldots,n)$, then it can be shown that, on taking $P(x_{i+1},\ldots,x_n)$ as basis, this module has in each case only finitely many elementary divisors distinct from the unit; this is a finiteness condition following from the existence of an ideal basis. The product over these finitely many elementary divisors gives the individual norms; the norm $N(\mathfrak a)$ is defined as the product of all individual norms. The largest primary factors $Q^{(i)}(x_i,x_{i+1},\ldots,x_n)$ of $N(\mathfrak a)$ correspond one-to-one to the associated prime ideals $\mathfrak p$ of $\mathfrak a$, in such a way that $Q^{(i)}$ becomes a power of $P^{(i)}$, where $P^{(i)}$ (or, if necessary, a power of $P^{(i)}$) denotes the norm $N(\mathfrak p)$ of $\mathfrak p$. Thus there arises a decomposition
\[
 N(\mathfrak a)=N(\mathfrak p_1)^{\varrho_1}\cdots N(\mathfrak p_r)^{\varrho_r}=Q_1\cdots Q_r,
\]
-- where if necessary $N(\mathfrak p)$ is to be replaced by the highest elementary divisor $E(\mathfrak p)=P^{(i)}$ -- and thus, by I, II, III, the decomposition of the norm into linear factors has been effected for all zeros of $\mathfrak a$. The multiplicity of the zeros corresponding to the individual prime ideal $\mathfrak p$ of $\mathfrak a$ can be regarded as the degree of $Q^{(i)}$, which can be interpreted as the number of residue classes linearly independent over $P(x_{i+1},\ldots,x_n)$. More precisely, these are the residue classes of certain isolated components of the decomposition uniquely determined by $\mathfrak p$ and by the prime ideals of higher dimension, namely the residue classes of the ground ideal of the $(i-1)$-st stage modulo the component, determined by $\mathfrak p$, of the ground ideal of the $i$-th stage; the latter component has $\mathfrak p$ and all prime ideals of higher dimension as associated prime ideals, while the former has only the prime ideals of higher dimension.

A complete \emph{insertion of the case of one indeterminate} is obtained when one passes here too to the principal ideal. The decomposition given under 1. then splits -- according to whether $a(x)$ is regarded as basis polynomial or norm -- into two separate decompositions: into the representation of the principal ideal as the least common multiple of largest primary components according to I, which here amounts to a representation as a product of powers of prime ideals; or into the decomposition of the norm as a product of powers of the norms of the individual prime ideals according to IV. In the construction of the zero field according to 2., what is in question in reality is the zero field of the prime ideal defined by $p(x)$, while in 3. the norm is decomposed into linear factors. The definition of multiplicity according to 4. is subordinated to that given in IV, since the ground ideal of the zeroth stage is equal to the unit ideal, while for one indeterminate the ground ideal of the first stage already equals the ideal itself.

The \emph{insertion into the usual elimination theory} is given by the fact that the norm of the ideal has the meaning of a resultant. In this interpretation, the old problem of elimination theory thus becomes a part of the ideal theory of the polynomial domain, namely that part which, beside the decomposition theorems for ideals, also treats the decomposition of their norms and the connection between the two decompositions.

\begin{center}
(Received 19 October 1923.)
\end{center}

\end{document}
